% GENERATED FILE. Edit canonical lessons, then run npm run generate:books.
% canonical-source-hash: fnv1a64:84d32e8b4bc7223d
% canonical-lessons: GU-C74-expensive, GU-C74-near, GU-C74-far, GU-C74-up, GU-C74-down

\chapter{Dear, Near, Far, Up and Down}
\label{ch:gu-dear-near-far-up-and-down}
\hlchaptermodality{74}

\begin{chapteropening}
\textbf{By the end of this chapter:} \emph{I can say expensive, near, far, up and down.}

Complete the last lesson of chapter 74: I can say expensive, near, far, up and down.
\end{chapteropening}

\section[\texorpdfstring{\gu{મોંઘું} (moṅghũ)}{મોંઘું (moṅghũ)}]{\gu{મોંઘું} (moṅghũ) — expensive}

\label{lesson:GU-C74-expensive}



\begin{quote}
Before the new one: say the Gujarati for hot, then the Gujarati for cheap.
\end{quote}

\subsection*{You'll want to know: \gu{મોંઘું}}
\textbf{\gu{મોંઘું}} (\emph{moṅghũ}) — \textquotedblleft{}expensive\textquotedblright{}. The opposite of \textbf{\gu{સસ્તું}}.

\begin{cousinweb}
\textbf{\gu{મોંઘું}} (\emph{moṅghũ}) is a tadbhava of Sanskrit \emph{mahārgha}, \textquotedblleft{}high-priced\textquotedblright{}, from \emph{mahā}, \textquotedblleft{}great\textquotedblright{}, and \emph{argha}, \textquotedblleft{}price\textquotedblright{} (CDIAL 9954).
\end{cousinweb}

\begin{grammarlens}[title={What you've built}]
One more word for dear, near, far, up and down.
\end{grammarlens}

\subsection*{Your turn}
\begin{itemize}
\raggedright
  \item \emph{Say it:} \emph{moṅghũ}
  \item \emph{Say it:} \emph{moṅghũ}, once more
  \item \emph{Say it:} say \emph{moṅghũ}
  \item \emph{Recall:} say the Gujarati for hot, then the Gujarati for cheap, then say \emph{moṅghũ} again
\end{itemize}

\begin{tcolorbox}[breakable,colback=teal!4,colframe=teal!35!black,title={Before you move on}]
What does \gu{મોંઘું} mean? (\textquotedblleft{}Expensive\textquotedblright{}.) Say it once more.
\end{tcolorbox}
\section[\texorpdfstring{\gu{નજીક} (najīk)}{નજીક (najīk)}]{\gu{નજીક} (najīk) — near}

\label{lesson:GU-C74-near}



\begin{quote}
Before the new one: say the Gujarati for cheap, then the Gujarati for expensive.
\end{quote}

\subsection*{You'll want to know: \gu{નજીક}}
\textbf{\gu{નજીક}} (\emph{najīk}) — \textquotedblleft{}near\textquotedblright{}.

\begin{cousinweb}
\textbf{\gu{નજીક}} (\emph{najīk}) is a loan, usually traced to Persian \emph{nazdīk}, \textquotedblleft{}near\textquotedblright{}. Gujarati writes the \emph{z} as \emph{j} and drops the \emph{d}.
\end{cousinweb}

\begin{grammarlens}[title={What you've built}]
One more word for dear, near, far, up and down.
\end{grammarlens}

\subsection*{Your turn}
\begin{itemize}
\raggedright
  \item \emph{Say it:} \emph{najīk}
  \item \emph{Say it:} \emph{najīk}, once more
  \item \emph{Say it:} say \emph{najīk}
  \item \emph{Recall:} say the Gujarati for cheap, then the Gujarati for expensive, then say \emph{najīk} again
\end{itemize}

\begin{tcolorbox}[breakable,colback=teal!4,colframe=teal!35!black,title={Before you move on}]
What does \gu{નજીક} mean? (\textquotedblleft{}Near\textquotedblright{}.) Say it once more.
\end{tcolorbox}
\section[\texorpdfstring{\gu{દૂર} (dūr)}{દૂર (dūr)}]{\gu{દૂર} (dūr) — far}

\label{lesson:GU-C74-far}



\begin{quote}
Before the new one: say the Gujarati for expensive, then the Gujarati for near.
\end{quote}

\subsection*{You'll want to know: \gu{દૂર}}
\textbf{\gu{દૂર}} (\emph{dūr}) — \textquotedblleft{}far\textquotedblright{}.

\begin{cousinweb}
\textbf{\gu{દૂર}} (\emph{dūr}) is a tadbhava, a word inherited from Sanskrit and reshaped by sound change. It comes from \emph{dūra}, \textquotedblleft{}distant\textquotedblright{} (CDIAL 6495).
\end{cousinweb}

\begin{grammarlens}[title={What you've built}]
One more word for dear, near, far, up and down.
\end{grammarlens}

\subsection*{Your turn}
\begin{itemize}
\raggedright
  \item \emph{Say it:} \emph{dūr}
  \item \emph{Say it:} \emph{dūr}, once more
  \item \emph{Say it:} say \emph{dūr}
  \item \emph{Recall:} say the Gujarati for expensive, then the Gujarati for near, then say \emph{dūr} again
\end{itemize}

\begin{tcolorbox}[breakable,colback=teal!4,colframe=teal!35!black,title={Before you move on}]
What does \gu{દૂર} mean? (\textquotedblleft{}Far\textquotedblright{}.) Say it once more.
\end{tcolorbox}
\section[\texorpdfstring{\gu{ઉપર} (upar)}{ઉપર (upar)}]{\gu{ઉપર} (upar) — up, above}

\label{lesson:GU-C74-up}



\begin{quote}
Before the new one: say the Gujarati for near, then the Gujarati for far.
\end{quote}

\subsection*{You'll want to know: \gu{ઉપર}}
\textbf{\gu{ઉપર}} (\emph{upar}) — \textquotedblleft{}up, above\textquotedblright{}.

\begin{cousinweb}
\textbf{\gu{ઉપર}} (\emph{upar}) is a tadbhava, a word inherited and reshaped by sound change. It comes from \emph{*uppari}, \textquotedblleft{}above\textquotedblright{}, a form CDIAL reconstructs from its descendants (CDIAL 2333). It is Sanskrit \emph{upari}, \textquotedblleft{}above\textquotedblright{}, with the consonant doubled for emphasis.
\end{cousinweb}

\begin{grammarlens}[title={What you've built}]
One more word for dear, near, far, up and down.
\end{grammarlens}

\subsection*{Your turn}
\begin{itemize}
\raggedright
  \item \emph{Say it:} \emph{upar}
  \item \emph{Say it:} \emph{upar}, once more
  \item \emph{Say it:} say \emph{upar}
  \item \emph{Recall:} say the Gujarati for near, then the Gujarati for far, then say \emph{upar} again
\end{itemize}

\begin{tcolorbox}[breakable,colback=teal!4,colframe=teal!35!black,title={Before you move on}]
What does \gu{ઉપર} mean? (\textquotedblleft{}Up, above\textquotedblright{}.) Say it once more.
\end{tcolorbox}
\section[\texorpdfstring{\gu{નીચે} (nīce)}{નીચે (nīce)}]{\gu{નીચે} (nīce) — down, below}

\label{lesson:GU-C74-down}



\begin{quote}
Before the new one: say the Gujarati for far, then the Gujarati for up.
\end{quote}

\subsection*{You'll want to know: \gu{નીચે}}
\textbf{\gu{નીચે}} (\emph{nīce}) — \textquotedblleft{}down, below\textquotedblright{}.

\begin{cousinweb}
\textbf{\gu{નીચે}} (\emph{nīce}) is the adverb of \emph{nīcũ}, \textquotedblleft{}low\textquotedblright{}, which CDIAL lists under Sanskrit \emph{nīca}, \textquotedblleft{}low\textquotedblright{} (CDIAL 7540). CDIAL explains the form as shaped by its opposite, \emph{ucca}, \textquotedblleft{}high\textquotedblright{}.
\end{cousinweb}

\begin{grammarlens}[title={What you've built}]
One more word for dear, near, far, up and down.
\end{grammarlens}

\subsection*{Your turn}
\begin{itemize}
\raggedright
  \item \emph{Say it:} \emph{nīce}
  \item \emph{Say it:} \emph{nīce}, once more
  \item \emph{Say it:} say \emph{nīce}
  \item \emph{Recall:} say the Gujarati for far, then the Gujarati for up, then say \emph{nīce} again
\end{itemize}

\begin{tcolorbox}[breakable,colback=teal!4,colframe=teal!35!black,title={Before you move on}]
What does \gu{નીચે} mean? (\textquotedblleft{}Down, below\textquotedblright{}.) Say it once more.
\end{tcolorbox}
